% !TeX root = ../main.tex
本节用冻结的 WOE Logistic 概率和未取整评分，在验证集选定规则后，以同一批
30,000 条测试记录评价。误报和漏报成本、审批坏率目标均为教学情景设定。

\subsection{代价敏感的风险标记}
设客户的坏样本概率为 $p$。标记好客户产生误报成本 $C_{\mathrm{FP}}=1$，
放过坏客户产生漏报成本 $C_{\mathrm{FN}}=5$。在概率可信且成本适用于每位客户时，
标记的条件期望成本为 $1-p$，放过的条件期望成本为 $5p$。取等号时标记，得到
\[
  1-p\leq 5p \quad\Longleftrightarrow\quad
  p\geq t_{\mathrm{theory}}=\frac{1}{1+5}=\frac16 .
\]
以冻结模型的预测概率 $\widehat p$ 执行时，由式~\eqref{eq:score-probability} 得到等价分数规则
\[
  \widehat p\geq t\quad\Longleftrightarrow\quad
  S\leq s(t),\qquad s(t)=A+B\ln\frac{1-t}{t}.
\]

验证前固定概率网格 $0.01,0.02,\ldots,0.99$，另纳入 $1/6$。以
$(\mathrm{FP}+5\,\mathrm{FN})/n$ 最小为准，同成本时取较高阈值以减少标记量。
验证集最优阈值为 $0.13$，均成本为 $0.21060$，对应 $S\leq 569.891525\ldots$。
同一验证集上，$0.5$ 和 $1/6$ 的均成本分别为 $0.28133$ 和 $0.21730$。
冻结 $0.13$ 后，按表~\ref{tab:decision-flag-results} 在测试集比较三种规则。

\begin{table}[H]
  \centering
  \caption{三种风险标记规则的测试集比较}
  \label{tab:decision-flag-results}
  \small
  \setlength{\tabcolsep}{3.5pt}
  \begin{tabular}{lrrrrrrrrr}
\toprule
规则 & 分数上限 & TN & FP & FN & TP & 标记率 & 召回率 & 精确率 & 均成本 \\
\midrule
基准 0.5 & 515.041 & 27,751 & 244 & 1,641 & 364 & 2.03\% & 18.15\% & 59.87\% & 0.28163 \\
理论 $1/6$ & 561.480 & 26,434 & 1,561 & 1,041 & 964 & 8.42\% & 48.08\% & 38.18\% & 0.22553 \\
验证 0.13 & 569.892 & 25,830 & 2,165 & 905 & 1,100 & 10.88\% & 54.86\% & 33.69\% & 0.22300 \\
\bottomrule
\end{tabular}


  \smallskip
\end{table}

表~\ref{tab:decision-confusion-matrices} 按实际类别展开三种规则的混淆矩阵。坏样本为正类，
好样本行中的标记人数为误报（FP），坏样本行中的未标记人数为漏报（FN）。

\begin{table}[H]
  \centering
  \caption{三种风险标记规则的测试集混淆矩阵}
  \label{tab:decision-confusion-matrices}
  \small
  \begin{tabular}{lrrrrrr}
    \toprule
    & \multicolumn{2}{c}{基准 \(t=0.5\)}
    & \multicolumn{2}{c}{理论 \(t=1/6\)}
    & \multicolumn{2}{c}{验证选定 \(t=0.13\)} \\
    \cmidrule(lr){2-3}\cmidrule(lr){4-5}\cmidrule(lr){6-7}
    实际类别 & 未标记 & 标记 & 未标记 & 标记 & 未标记 & 标记 \\
    \midrule
    好样本 & 27,751 & 244 & 26,434 & 1,561 & 25,830 & 2,165 \\
    坏样本 & 1,641 & 364 & 1,041 & 964 & 905 & 1,100 \\
    \bottomrule
  \end{tabular}
\end{table}

$0.5$ 只标记 608 户，精确率为 59.87\%，但漏掉 1,641 条坏样本，均成本最高。
由 $1/6$ 降至 $0.13$ 后，标记量由 2,525 户增至 3,265 户，少漏 136 条坏样本，多误报 604 户。
少漏带来的成本节省为 $136\times5=680$，扣除新增误报成本后净省 76，均成本由 0.22553 降至 0.22300。
验证阈值在三种规则中成本最低，同时需要处理测试集约 10.88\% 的客户，精确率降至 33.69\%。

\Needspace{23\baselineskip}
\subsection{坏率约束下的模拟审批}
给定整数分数门槛 $c$，接受所有未取整评分 $S_i\geq c$ 的客户，同分者使用同一规则。
设验证集接受人数为 $N_c$、其中坏样本数为 $D_c$、验证集总人数为 $n_{\mathrm{val}}$。
在事先固定的 $c=300,301,\ldots,900$ 上求
\[
  \max_c\ \frac{N_c}{n_{\mathrm{val}}}
  \qquad\text{满足}\qquad
  \frac{D_c}{N_c}\leq 5\%,\qquad
  \frac{N_c}{n_{\mathrm{val}}}\geq 10\%.
\]

同接受人数时取较高门槛，若无可行门槛则报告不可行。
验证集在 530 分时接受 28,973 户，坏率为 5.03\%，超出约束。
遍历网格后，531 分是可行门槛中接受人数最多的方案，接受 28,928 户，坏率为 4.97\%。
测试阶段沿用这一门槛，结果见表~\ref{tab:decision-approval-results}。

\begin{table}[H]
  \centering
  \caption{冻结审批门槛在验证集与测试集上的表现}
  \label{tab:decision-approval-results}
  \small
  \setlength{\tabcolsep}{5pt}
  \begin{tabular}{lrrrrrrr}
\toprule
样本 & 门槛 & 接受人数 & 接受率 & 坏样本 & 平均预测 & 实际坏率 & 95\% 区间 \\
\midrule
验证集 & 531 & 28,928 & 96.43\% & 1,437 & 4.72\% & 4.97\% & -- \\
测试集 & 531 & 28,974 & 96.58\% & 1,465 & 4.79\% & 5.06\% & ［4.81\%, 5.31\%］ \\
\bottomrule
\end{tabular}


  \smallskip
  \parbox{0.96\linewidth}{\small 区间为测试集接受样本实际坏率的记录层面 Wilson 95\% 二项比例区间。}
\end{table}

测试集接受率为 96.58\%，接受者平均预测为 4.79\%，实际坏率为 5.06\%。
点估计超过 5\% 目标约 0.06 个百分点，区间 $[4.81\%,5.31\%]$ 跨越 5\%。
据此记录测试点估计超标，保留冻结门槛 531 分。区间表明，接受者坏率在 5\% 附近仍有抽样不确定性。

按本报告的评分刻度，600 分对应好坏几率 $19:1$，预测坏样本概率为 $1/(19+1)=5\%$。
对一位恰好得 600 分的客户，这表示模型根据评分时已知的信息，估计其未来两年内发生严重逾期或更严重情形的概率为 5\%。
若规定 $S\geq 600$ 时接受，每位接受者的预测概率都不超过 5\%，客群的平均预测概率也不超过 5\%。
实际坏样本率则要等两年表现窗口结束后，用发生目标事件的人数除以接受人数计算。
受模型校准偏差和样本波动影响，实际坏样本率可能超过 5\%。本方案约束的是验证集的客群实际坏样本率，
因此允许部分预测概率高于 5\% 的客户被接受。

\paragraph{两种规则的行动目标}
风险标记用于安排复核，在给定误报和漏报代价下比较标记量、检出量与成本。
模拟审批则在客群坏率约束下扩大接受范围。例如，未取整分数在 $531$ 到 $569.891525\ldots$ 之间的客户，
可同时满足接受与风险标记条件。现有数据支持比较复核容量和模拟接受风险。
若要评价实际授信收益或优化审批政策，还需结合贷款收益、违约损失率与拒绝样本的后续表现。
